\section{Work with one cumulative source}
The cumulative starter is \path{packages/027-rsa-arithmetic/starter/main.asm}. Keep one working
copy and replace the marked region bodies as you reach their steps. BEGIN and END
comments identify the replacement boundaries. Some bodies contain loop labels; copy
those labels along with their instructions. The supplied caller, frames and restores
remain in place. There is no separate finished template to switch to halfway through.

The default caller selects the gcd. Prepare other observations from the
\path{packages/027-rsa-arithmetic} directory. For example, invoke \texttt{python3 prepare.py s13-1}
with \texttt{--source} naming your cumulative working file and \texttt{--output} naming
a new observation file. The helper initializes only the declared data arenas and mode.
It retains your authored code and refuses to overwrite an existing output. Use the
cumulative version through the named edit, since later edits can change a temporary
observation into a loop that continues farther.

Set a breakpoint at the recipe's named checkpoint and inspect it before continuing.
A checkpoint named \texttt{function\_return} means the corresponding
\texttt{gcd64\_return}, \texttt{pow32\_return} or other entry-specific return label when
using the supplied caller. Direct function checks use \texttt{function\_return} itself.
The caller stores status after that return boundary. For gcd it stores \texttt{v1},
since \texttt{v0} contains the numerical gcd; other entries use \texttt{v0} for status.

Development status three means that a region is unfinished. Intermediate bytes are
observable for learning but are not promised as a successful public result. Nested
calls must check status before consuming temporary data, then return through their own
matching restore on failure. The completed success/error contracts remain the ones in
the interface table. Guided observations describe the illustrated implementation;
free-writing checks accept another algorithm that satisfies those contracts.
